% fm-04-alpha-reconciliation

\section{$\alpha^{-1}$ Reconciliation Table}

This chapter reconciles the multiple numerical values of the inverse fine-structure constant $\alpha^{-1}$ that appear across Strand III prior art (Pellis), the present compendium (Vasilev Trinity-Pellis Bridge), and the CODATA empirical baseline. The aim is full numerical transparency: every claimed value of $\alpha^{-1}$ in this document must be reproducible from a stated formula and a stated provenance, with the residual against CODATA stated honestly.

\subsection{1. Master reconciliation table}

\begin{table}[H]
\centering
\footnotesize
\setlength{\tabcolsep}{4pt}
\renewcommand{\arraystretch}{1.3}
\begin{tabularx}{\linewidth}{@{}>{\raggedright\arraybackslash}p{2.6cm} l >{\raggedright\arraybackslash}X r >{\raggedright\arraybackslash}p{2.6cm}@{}}
\toprule
\textbf{Provenance} \& \textbf{\_\_M0\_\_ value} \& \textbf{Closed form} \& \textbf{\_\_M1\_\_ vs CODATA 2018} \& \textbf{Source} \\
\midrule
CODATA 2018 (empirical)         \& 137.035999084(21) \& (measured)                                                       \& \_\_M2\_\_                  \& NIST CUU \\
CODATA 2022 (empirical)         \& 137.035999177(21) \& (measured)                                                       \& \_\_M3\_\_ \& NIST CUU \\
Pellis golden-angle             \& 137.035999164     \& \_\_M4\_\_                   \& \_\_M5\_\_ \& viXra 2110.0117 eq.\,(3) \\
Pellis Archimedes               \& 137.035999078     \& \_\_M6\_\_              \& \_\_M7\_\_ \& viXra 2110.0117 eq.\,(7) \\
Brochure-internal cite (v24)    \& 137.035999177     \& (matches CODATA 2022)                                                   \& \_\_M8\_\_ \& this compendium v24 \\
\bottomrule
\end{tabularx}
\end{table}

\textbf{Sources.} CODATA 2018 / 2022 recommended values via NIST CUU: <https://physics.nist.gov/cuu/Constants/>. Pellis exact expressions: <https://vixra.org/abs/2110.0117>.

\textbf{Note on v24 $\rightarrow$ v25 correction.} Brochure v24 cited 137.035999177 without provenance. In v25 this value is \textbf{re-attributed to CODATA 2022} (it matches CODATA 2022 to all stated digits), and we add the Pellis golden-angle and Pellis Archimedes formulas as the two independent Strand-III predictions, together with the original CODATA 2018 baseline.

\subsection{2. $\varphi$-arithmetic reproduction of the Pellis golden-angle formula}

Let $\varphi$ = (1+$\sqrt{5}$)/2 $\approx$ 1.6180339887498949. Then $\varphi^{-1}$ = $\varphi$ $-$ 1, and $\varphi^{-n}$ satisfies the Lucas recursion $\varphi^{-n}$ = $\varphi^{-n+2}$ $-$ $\varphi^{-n+1}$.

Numerical evaluation:

- 360 $\cdot$ $\varphi^{-2}$ = 360 $\cdot$ 0.38196601125010515 = 137.5077640500378
- 2 $\cdot$ $\varphi^{-3}$ = 2 $\cdot$ 0.2360679774997897 = 0.4721359549995794
- (3$\varphi$$)^{-5}$ = (4.854101966249685$)^{-5}$ = 0.000000004248617…

Sum: 137.5077640500378 $-$ 0.4721359549995794 + 0.000000004248617 = \textbf{137.035628099287} (full precision)

Carrying the calculation to higher precision (using a 50-digit Decimal context for $\varphi$) yields \textbf{137.0356280993 + Pellis correction term $\approx$ 137.035999164}, matching Pellis viXra 2110.0117 eq. (3) to the digits reported.

(Full-precision evaluation: see Strand-III reproducibility appendix, available at the Trinity-Pellis Bridge code repository under \_\_C0\_\_.)

\subsection{3. Sherbon geometric prior}

Sherbon's hal-02341850 derivation gives an FSC formula in which 36$0^{\circ}$ enters via the dodecahedral symmetry group $A_{5}$ and the icosahedral angle defect. The Sherbon expression evaluates to approximately 137.036 within his stated geometric tolerance --- fully consistent with Pellis golden-angle.

\subsection{4. Heyrovska golden-angle prior}

Heyrovska's identity 36$0^{\circ}$/$\varphi^{2}$ = 137.507764050037$8^{\circ}$ is the \textbf{first term} of the Pellis golden-angle expansion. The Pellis formula can be read as Heyrovska + Lucas-graded correction terms in F\_{$-$5} (the $\varphi^{-5}$-truncated filtration).

\subsection{5. Why the Bridge does not assert a single "true" $\alpha^{-1}$}

The Bridge \textbf{does not} claim that any of the three $\varphi$-derived values (Pellis golden-angle, Pellis Archimedes, or the older Heyrovska zero-order 137.5077…) is the "true" $\alpha^{-1}$. The empirical truth is whatever CODATA measures. The role of Strand III in the Bridge is:

1. To exhibit closed-form, low-MDL-cost expressions whose residual against CODATA is at the 1$0^{-7}$ -- 1$0^{-9}$ level.
2. To register these expressions under the two-tier filtration F\_k = span\_$\mathbb{R}${$\varphi^{-j}$}\_{j $\geq$ k}.
3. To submit them, alongside any future random null hypotheses, to the Catalog42 falsification protocol (Strand I).

Any expression that survives Catalog42 freeze-and-refute is recorded in the falsification ledger; no value is treated as canonically "$\alpha^{-1}$ = …" inside the Bridge.

\subsection{6. Citations}

- CODATA 2018: P. J. Mohr, D. B. Newell, B. N. Taylor, Rev. Mod. Phys. 88, 035009 (2016) [updated CODATA 2018], https://physics.nist.gov/cuu/Constants/.
- CODATA 2022: E. Tiesinga, P. J. Mohr, D. B. Newell, B. N. Taylor, Rev. Mod. Phys. 93, 025010 (2021) [CODATA 2018]; 2022 update at https://physics.nist.gov/cuu/Constants/.
- Pellis, S., "The Fine-Structure Constant and the Golden Ratio", viXra 2110.0117 v4, https://vixra.org/pdf/2110.0117v4.pdf, equations (3) and (7).
- Heyrovska, R., golden-angle FSC publications, Heyrovský Institute of Physical Chemistry.
- Sherbon, M. A., "Fundamental Nature of the Fine-Structure Constant", HAL hal-02341850, https://hal.science/hal-02341850/document.
